\documentclass[11pt]{article}
\usepackage[margin=1in]{geometry}
\usepackage{amsmath, amssymb}
\usepackage{enumitem}
\usepackage{tcolorbox}
\usepackage{xcolor}

\title{\textbf{Homework 6, Question 4c Explanation} \\ \large Poisson Process with Random Classification}
\author{}
\date{}

\begin{document}

\maketitle

\section{Problem Statement}

\textbf{Given:}
\begin{itemize}[noitemsep]
    \item $\{N_t\}$ is a Poisson process with rate $\lambda$
    \item Each arrival is \textbf{independently} classified as:
    \begin{itemize}
        \item Woman with probability $\frac{1}{2}$
        \item Man with probability $\frac{1}{2}$
    \end{itemize}
    \item The genders are independent of each other and of the arrival times
    \item $M_t$ = number of women that arrived by time $t$
\end{itemize}

\textbf{Find:} $\mathbb{E}(N_t \mid M_t = m)$

\section{Key Concept: Poisson Thinning}

\begin{tcolorbox}[colback=blue!5!white,colframe=blue!75!black,title=Poisson Thinning Theorem]
When you ``thin'' a Poisson process with rate $\lambda$ by randomly and independently classifying each arrival with probability $p$:
\begin{itemize}
    \item The ``successes'' form a Poisson process with rate $p\lambda$
    \item The ``failures'' form a Poisson process with rate $(1-p)\lambda$
    \item \textbf{These two processes are independent}
\end{itemize}
\end{tcolorbox}

\subsection{Applying Thinning to Our Problem}

In this problem, we're classifying each arrival as woman (probability $\frac{1}{2}$) or man (probability $\frac{1}{2}$):

\begin{align*}
M_t &= \text{number of women by time } t \sim \text{Poisson}\left(\frac{\lambda t}{2}\right) \\
W_t &= \text{number of men by time } t \sim \text{Poisson}\left(\frac{\lambda t}{2}\right) \\
N_t &= M_t + W_t \quad \text{(total arrivals)}
\end{align*}

\begin{tcolorbox}[colback=red!5!white,colframe=red!75!black]
\textbf{Crucial Fact:} $M_t$ and $W_t$ are \textbf{independent} random variables.
\end{tcolorbox}

\section{Solution}

We want to find $\mathbb{E}(N_t \mid M_t = m)$.

Since $N_t = M_t + W_t$, we can write:

\begin{align}
\mathbb{E}(N_t \mid M_t = m) &= \mathbb{E}(M_t + W_t \mid M_t = m) \notag \\
&= \mathbb{E}(M_t \mid M_t = m) + \mathbb{E}(W_t \mid M_t = m) \label{eq:linearity}
\end{align}

\subsection{Step 1: Evaluate $\mathbb{E}(M_t \mid M_t = m)$}

This is straightforward:
\[
\mathbb{E}(M_t \mid M_t = m) = m
\]

If we know $M_t = m$, then the expected value of $M_t$ is simply $m$.

\subsection{Step 2: Evaluate $\mathbb{E}(W_t \mid M_t = m)$}

\begin{tcolorbox}[colback=green!5!white,colframe=green!75!black,title=Key Step: Using Independence]
Since $M_t$ and $W_t$ are \textbf{independent}, conditioning on $M_t = m$ does not affect the distribution of $W_t$:
\[
\mathbb{E}(W_t \mid M_t = m) = \mathbb{E}(W_t)
\]
\end{tcolorbox}

Therefore:
\[
\mathbb{E}(W_t) = \frac{\lambda t}{2}
\]

This is just the mean of a Poisson random variable with parameter $\frac{\lambda t}{2}$.

\subsection{Step 3: Combine the Results}

Substituting back into equation \eqref{eq:linearity}:

\begin{align*}
\mathbb{E}(N_t \mid M_t = m) &= m + \mathbb{E}(W_t) \\
&= m + \frac{\lambda t}{2}
\end{align*}

\begin{tcolorbox}[colback=yellow!10!white,colframe=orange!75!black,title=Final Answer]
\[
\boxed{\mathbb{E}(N_t \mid M_t = m) = m + \frac{\lambda t}{2}}
\]
\end{tcolorbox}

Alternatively, since $\mathbb{E}(M_t) = \frac{\lambda t}{2}$, we can write:
\[
\boxed{\mathbb{E}(N_t \mid M_t = m) = m + \mathbb{E}(M_t)}
\]

\section{Intuition}

Let's interpret what this result means:

\begin{center}
\textit{``Given that we observed $m$ women, how many total people do we expect?''}
\end{center}

\textbf{Answer:}
\begin{itemize}
    \item We know \textbf{for certain} there were $m$ women
    \item The number of men is \textbf{independent} of the number of women (because each person's gender is determined independently)
    \item So we expect the \textbf{average} number of men: $\mathbb{E}(W_t) = \frac{\lambda t}{2}$
    \item Total expected = $m + \frac{\lambda t}{2}$
\end{itemize}

\subsection{Why Independence Matters}

The key insight is that \textbf{independence} means:
\begin{center}
\textit{Knowing $M_t = m$ gives us no information about $W_t$}
\end{center}

Therefore:
\[
\mathbb{E}(W_t \mid M_t = m) = \mathbb{E}(W_t)
\]

If $M_t$ and $W_t$ were \textit{not} independent, we couldn't make this simplification!

\section{Numerical Example}

\textbf{Example:} Suppose $\lambda = 4$ arrivals per hour, $t = 2$ hours, and we observe $m = 3$ women.

\textbf{What is the expected total number of arrivals?}

\textbf{Solution:}
\begin{align*}
\mathbb{E}(N_t \mid M_t = 3) &= 3 + \frac{\lambda t}{2} \\
&= 3 + \frac{4 \times 2}{2} \\
&= 3 + 4 \\
&= 7 \text{ arrivals}
\end{align*}

\textbf{Interpretation:}
\begin{itemize}
    \item We know there were 3 women
    \item On average, we expect $\frac{4 \times 2}{2} = 4$ men
    \item Total expected: $3 + 4 = 7$ people
\end{itemize}

\section{Common Mistakes to Avoid}

\subsection{Mistake 1: Forgetting Independence}

\textcolor{red}{\textbf{Wrong:}} ``Since we saw more women than expected, there must be fewer men''

\textcolor{green!60!black}{\textbf{Correct:}} The number of men is independent of the number of women, so observing $M_t = m$ doesn't change our expectation for $W_t$.

\subsection{Mistake 2: Using $\mathbb{E}(N_t)$ Instead}

\textcolor{red}{\textbf{Wrong:}} $\mathbb{E}(N_t \mid M_t = m) = \mathbb{E}(N_t) = \lambda t$

\textcolor{green!60!black}{\textbf{Correct:}} Conditioning on $M_t = m$ changes the expectation. We need to account for the fact that we \textit{know} there are $m$ women.

\subsection{Mistake 3: Double Counting}

\textcolor{red}{\textbf{Wrong:}} $\mathbb{E}(N_t \mid M_t = m) = m + \lambda t$ (counting $m$ women plus all expected arrivals)

\textcolor{green!60!black}{\textbf{Correct:}} We add $m$ (known women) plus $\frac{\lambda t}{2}$ (expected men only), not all arrivals.

\section{Related Facts}

\subsection{Conditional Distribution}

Not only is the conditional \textit{expectation} straightforward, but the full conditional \textit{distribution} is also simple:

\[
(N_t \mid M_t = m) \sim m + \text{Poisson}\left(\frac{\lambda t}{2}\right)
\]

In other words, $N_t$ given $M_t = m$ is distributed as $m$ (deterministic) plus an independent Poisson random variable with rate $\frac{\lambda t}{2}$.

\subsection{Variance}

What about $\text{Var}(N_t \mid M_t = m)$?

\begin{align*}
\text{Var}(N_t \mid M_t = m) &= \text{Var}(M_t + W_t \mid M_t = m) \\
&= \text{Var}(m + W_t \mid M_t = m) \\
&= \text{Var}(W_t \mid M_t = m) \\
&= \text{Var}(W_t) \quad \text{(by independence)} \\
&= \frac{\lambda t}{2}
\end{align*}

So the conditional variance equals the variance of $W_t$, which makes sense since $m$ is constant when conditioning on $M_t = m$.

\section{Summary}

\begin{tcolorbox}[colback=gray!5!white,colframe=gray!75!black,title=Key Takeaways]
\begin{enumerate}
    \item \textbf{Poisson thinning} splits a Poisson process into independent components
    \item When random variables are \textbf{independent}, conditioning on one doesn't affect the other
    \item The answer $m + \frac{\lambda t}{2}$ = (known women) + (expected men)
    \item This generalizes: if probability of woman is $p$, then
    \[
    \mathbb{E}(N_t \mid M_t = m) = m + (1-p)\lambda t
    \]
\end{enumerate}
\end{tcolorbox}

\end{document}
